\begin{pastexam}[label={px:ntu-112-ee-ds-03}]{臺大電機所 112 學年度 第 3 題}
（是非；命題正確填 (A)，錯誤填 (B)；答對得 4 分，答錯倒扣 4 分，未填答不計分，倒扣至是非題總分 0 分為止）

Since a \textit{Deap} contains a subtree of \textit{min-heap} and a subtree of \textit{max-heap}, its \texttt{find()}, \texttt{deleteMin()}, and \texttt{deleteMax()} operations can all have \texttt{O(log n)} time complexities.
\end{pastexam}

\begin{pastsoln}
答案是 (B)，錯誤。

deap 的根不放元素，左子樹是最小堆積、右子樹是最大堆積，左邊每個位置不大於右邊的對應點（資結書
第~8.5.2~節）。 \texttt{deleteMin()}、\texttt{deleteMax()} 確實是 $O(\log n)$：洞沿著一邊往下走到底，再
把最後一個元素放回去、上浮，都只走一條路徑。

\texttt{find()} 不行。 找一個任意的鍵，堆積的順序只說「父節點不大於（不小於）子節點」，沒有告訴你目標
在左子樹還是右子樹，也不知道在最小堆積還是最大堆積那一邊。 最壞要看過所有元素，$\Theta(n)$。 例如在
最小堆積那一邊，所有葉節點之間沒有大小關係，要找的鍵可能在任何一片葉子上，而葉節點約有 $n/4$ 個。
所以三個操作「都」是 $O(\log n)$ 不成立。

若題目的 find 指的是找最小值和最大值，那是 $O(1)$（在 $A[2]$ 和 $A[3]$），但一般的 \texttt{find()} 是
搜尋任意鍵。
\end{pastsoln}
